\subsection{The cusp and tangent families together}
\label{sec:cusptangent}

Both remaining irreducible or line--conic cases have a unique
leading singular direction. If $F_2$ is nonzero there, a gradient
argument excludes purity. If it is zero and the direction is
isotropic, the frame test applies. We treat the nonisotropic
alternative with its full Gaussian
\[
 q=\tfrac12(Au^2+Bv^2+mT^2+2Cuv+2DuT+2EvT)
                  +\ell_u u+\ell_v v+\ell_TT+q_0,
 \qquad m\ne0.
\]
Translations and linear changes give the two exact normal forms
\begin{align}
 P_C&=T(u^2+\gamma v+h)+v^3+buv+eu+fv+k,
                       &&\hbox{cusp},\label{eq:cuspnormal}\\
 P_T&=T(u^2+gv+h)+(u+c)v^2+eu+fv+k,
                       &&\hbox{tangent line--conic}.\label{eq:tangentnormal}
\end{align}
Here is the leading-form reduction. At the unique singular
point, $F_3$ is $T$ times a rank-one binary quadratic plus a
binary cubic. Choose that quadratic as $u^2$. A linear shift
of $T$ removes the $u^3,u^2v$ terms, leaving
$Tu^2+a uv^2+b_0v^3$. If $b_0\ne0$, a shift of $v$ by a
multiple of $u$ removes $uv^2$, and rescaling gives the cusp.
If $b_0=0$, tangency to a smooth conic requires $a\ne0$;
rescaling gives the tangent form. The direction of $T$ is
unchanged, so its Gaussian square coefficient remains nonzero.
For the cusp, translate $u$ to remove $uT$, translate $T$ to
remove $u^2$, and translate $v$ to remove $v^2$. For the tangent
case use the same first two translations, then a $v$ translation
to remove $uv$. These operations leave the displayed arbitrary
coefficients and transform, rather than reset, the means and
cross terms of $q$.

\begin{lemma}\label{lem:cusptangentgate}
For both leading forms, $F_2$ nonzero at the singular direction
implies $\|\nabla P\|\ge c\|x\|$ outside a complex compact set.
In the nonisotropic forms above, each fixed nonzero phase
$tP-q$ has the same type of lower bound, locally uniformly in
$t$. Their full generic twisted ranks are six and five,
respectively.
\end{lemma}
\begin{proof}
Any failure of a linear gradient bound approaches the unique
leading singular direction, so $u,v=o(T)$. For the cusp,
the $u$ equation divided by $T$ makes $u$ bounded. The $T$
equation divided by $T$ then has the nonzero limit $2F_2(0,0,1)$.
For the tangent, the transverse equations give
$uv=O(T)$ and $2Tu+v^2=O(T)$. If $u$ is unbounded these imply
$u=O(|T|^{1/3})$; otherwise $u^2/T\to0$ directly. The $T$
equation again has the same nonzero limit. Apply the Gaussian
gradient criterion. Replacing the quadratic value by $-m/2$
gives the fixed-phase bounds for the normal forms. Compactness
of the coefficient estimates gives local parameter uniformity.

For clarity, the ranks do not come from an escape count. The
three quadratic gradient equations have total projective length
eight. Their only infinity point is the $T$ direction. In the
chart $T=1$ with homogenizing coordinate $z$, the cusp $u$
equation solves $u=O(z)$; the $T$ equation then gives $z=0$,
and the last equation is $3tv^2=0$. Its length is two. For
the tangent the $T$ equation solves $z=O(u^2)$, the $u$ equation
gives $u=-v^2/2+O(v^4)$, and the $v$ equation has leading term
$-tv^3$. Its length is three. Fixed-phase coercivity excludes
affine components of positive dimension. B\'ezout and Pham's
tame comparison therefore give ranks $8-2=6$ and $8-3=5$.
\end{proof}

Every escape in these two normal forms approaches the $T$
direction, where $q_2=m/2\ne0$. Proposition~\ref{prop:boundedregularrank}
therefore applies. We now list all possible bounded contributions;
the other exponential blocks require no explicit construction.

\subsubsection*{Isolated critical points on a smooth constraint}

Write $P=T A(u,v)+G(u,v)$. Where $A=0$ is smooth, take $s=A$
and a coordinate $y$ along that curve. The analytic shear
\[
 T'=T+\frac{G(y,s)-G(y,0)}s
\]
gives the exact identity $P=sT'+G(y,0)$. At an isolated critical
point this becomes $\alpha+sT'+y^{\mu+1}$. The already explicit
characters \eqref{eq:ordinarycharacters} provide all $\mu$
dimensions with positive original-marker order.

For the cusp, $\gamma\ne0$ gives a degree-six restriction to
$v=-(u^2+h)/\gamma$, so total isolated multiplicity is five.
If $\gamma=0,h\ne0$, the two lines $u=\pm\sqrt{-h}$ each
have a cubic restriction in $v$, giving multiplicity two each.
For the tangent, $g\ne0$ gives a degree-five restriction and
multiplicity four. If $g=0,h\ne0$, each line $u=u_i$ gives
\[
 (u_i+c)v^2+fv+eu_i+k.
\]
It has one Morse point when $u_i+c\ne0$, none when
$u_i+c=0,f\ne0$, and a critical parabola when both vanish.
The last case is included below. All repeated roots and equal
critical values in the isolated cases retain their Milnor
multiplicities; no Morse perturbation is assumed.

\subsubsection*{Four fixed models along the $T$ lines}

Here and below a scalar constant is harmless but retained in the
phase value. Set $\epsilon=1/t$. The critical equation in $T$
is
\begin{equation}\label{eq:cuspcriticalbound}
 u^2=\epsilon(mT+Du+Ev+\ell_T)
\end{equation}
whenever the constraint is the double line $u^2=0$.
Thus every bounded phase-critical branch has $u=O(\epsilon^{1/2})$.
This estimate will also verify that the following scalings lose
no bounded contribution.

First consider a critical $T$ line at $u=0,v=v_0$ whose transverse
quadratic part is $Tu^2+\lambda(v-v_0)^2+\beta u(v-v_0)$,
with $\lambda\ne0$. Scale $u=t^{-1/2}U$,
$v=v_0+t^{-1/2}V$, leaving $T$ fixed. The volume is $t^{-1}$
and the fixed phase is
\[
 H_0=TU^2+\lambda V^2+\beta UV-mT^2/2-\eta T-q_*.
\]
Complete $V$, then $T$. This is a quartic in $U$, with leading
coefficient $1/(2m)$, plus two nondegenerate Gaussians. Its
twisted rank and critical Milnor number are both three, including
all quartic collisions. Equation~\eqref{eq:cuspcriticalbound}
and the $v$ equation with nonzero linear coefficient $2\lambda$
give $v-v_0=O(t^{-1/2})$. Every bounded branch in this group has
bounded scaled variables. The model therefore supplies all three
local dimensions, with original-marker values in
$e^{\alpha t}t^{-1}\C[[t^{-1/2}]]$.

If instead $\lambda=0,\beta\ne0$, the same scaling gives
\[
 H_0=TU^2+\beta UV-mT^2/2-\eta T-q_*.
\]
The polynomial shear $V\mapsto V+TU/\beta$ removes $TU^2$.
This is a nondegenerate three-variable Gaussian, of rank one.
The $u$ equation with coefficient $\beta$ bounds $v-v_0$ by
$O(u)+O(\epsilon)$, so bounded branches are again captured.
It gives one dimension with the same positive volume order.

For the cusp's remaining line, $P=T u^2+v^3+k$, use
$u=t^{-1/2}U$, $v=t^{-1/3}V$. The fixed phase is
\[
 TU^2+V^3-mT^2/2-\ell_TT-q_0.
\]
Completing $T$ gives a quartic, a cubic and one Gaussian, of
total rank and Milnor number $3\cdot2=6$. The volume is
$t^{-5/6}$. Bounded critical branches satisfy $u=O(t^{-1/2})$
and, from $3v^2=\epsilon q_v$, $v=O(t^{-1/2})$, so their
scaled coordinates are bounded.

For the tangent's remaining line, $P=T u^2+uv^2+k$, use
$u=t^{-1/2}U$, $v=t^{-1/4}V$. Completing $T$ gives the fixed
two-variable phase
\[
 U^4/(2m)+UV^2-\ell_TU^2/m
\]
and one Gaussian. Its rank and critical Milnor number are five.
Here is a direct check: give $U,V$ weights $1/4,3/8$. The
leading derivatives are $2U^3/m+V^2$ and $2UV$, with sole
common zero the origin and quotient basis
$1,U,U^2,V,V^2$. They form a regular sequence. The associated
graded twisted de Rham complex is their Koszul complex,
concentrated in top cohomology; all derivative and lower-phase
terms have smaller weight. It follows that the same five
classes give the twisted rank. The ordinary gradient quotient
has the same weighted leading calculation, so its total
Milnor number is also five for every $\ell_T$. The volume is
$t^{-3/4}$. Equation~\eqref{eq:cuspcriticalbound} bounds $u$,
and $2Tu+v^2=\epsilon q_u$ gives $v=O(t^{-1/4})$. This
captures every bounded branch.

All remaining terms in these four scaled phases have positive
order in a common finite cover. Their coefficients are
polynomials at each order, and every fixed original marker has
nonnegative order. Formula~\eqref{eq:boundedfixedmodel} constructs
the stated complete local functionals. The Jacobian factors and
coordinate-velocity terms are retained, exactly as in its proof.

\subsubsection*{The critical parabola}

This occurs only for the tangent form, with $g=0$, $h\ne0$,
$c^2+h=0$, $f=0$, $c\ne0$. Put
\[
 x=u+c,\qquad s=(x-2c)T+v^2+e,\qquad\alpha=k-ec.
\]
Exactly $P=\alpha+xs$. Write $x=\rho X$, $s=\rho Y$,
$\rho=t^{-1/2}$. The volume is
\[
 \frac{\rho^2}{\rho X-2c}\,dX\,dv\,dY,
\]
and the fixed phase is $XY+\psi(v)$, where
\[
 \psi(v)=-q\!\left(-c,v,\frac{v^2+e}{2c}\right).
\]
Its quartic coefficient is $-m/(8c^2)\ne0$, so its fixed
twisted rank and critical Milnor number are three. At every
bounded phase-critical branch near this parabola, the exact
normal equations give $x=\epsilon q_s=O(\epsilon)$ and
$s=\epsilon q_x=O(\epsilon)$; the denominators stay away
from zero. Thus $X,Y$ are bounded. The scaled Jacobian is a
unit times $\rho^2$, and its expansion and that of the phase
have polynomial coefficients at each order. This gives all
three dimensions, with strictly positive original-marker order.
The other constraint line has its one isolated Morse point.

\subsubsection*{Exhaustive routing and completeness}

The remaining double-line cases require only the following two
tables. Their rows are disjoint. A missing common root means
there is no bounded contribution at all, not an omitted escape.

\begin{center}\small
\begin{tabular}{@{}p{0.36\textwidth}p{0.54\textwidth}@{}}
\toprule
Cusp with $\gamma=h=0$ & Bounded $P$-critical groups\\
\midrule
$f\ne0$ & For each root $3v_i^2+f=0$ with $bv_i+e=0$,
one line with $\lambda=3v_i$, $\beta=b$, rank three\\
$f=0$, $e\ne0$ & None, for every $b$\\
$f=e=0$, $b\ne0$ & One line with $\lambda=0$, $\beta=b$, rank one\\
$f=e=b=0$ & Quartic--cubic model, rank six\\
\bottomrule
\end{tabular}
\end{center}
The line values in the first row are
$\alpha_i=v_i^3+fv_i+k$; the other line values are $k$.

\begin{center}\small
\begin{tabular}{@{}p{0.36\textwidth}p{0.54\textwidth}@{}}
\toprule
Tangent with $g=h=0$ & Bounded $P$-critical groups\\
\midrule
$c\ne0$ & Put $v_0=-f/(2c)$. None if $e+v_0^2\ne0$;
otherwise one rank-three line with $\lambda=c$, $\beta=2v_0$,
value $k-f^2/(4c)$\\
$c=0$, $f\ne0$ & None, for every $e$\\
$c=f=0$, $e\ne0$ & Two lines $v_i^2=-e$, each with
$\lambda=0$, $\beta=2v_i$, rank one, value $k$\\
$c=f=e=0$ & Weighted quartic model, rank five, value $k$\\
\bottomrule
\end{tabular}
\end{center}

To verify these tables directly, the double constraint forces
$u=0$. The cusp equations reduce to $3v^2+f=0$, $bv+e=0$;
the tangent equations reduce to $2cv+f=0$, $v^2+e=0$.
The remaining $T$ coordinate is free. Together with the smooth
constraint and separated-line alternatives already proved,
this exhausts all coefficients of both normal forms.

For every fixed model the critical scheme is finite with the
stated total Milnor number. The bounds above show that every
bounded phase-critical branch has bounded scaled coordinates.
Conservation of local critical multiplicity thus gives exactly
the stated group counts. Proposition~\ref{prop:boundedregularrank}
shows that the groups at value zero account for the entire
regular rank; all escaping branches are excluded from it.

We also need independence on the original polynomial module.
Within a fixed group evaluate a polynomial basis of its fixed
quotient: its normalized leading matrix is invertible by
\eqref{eq:boundedfixedmodel}. Isolated groups use their finite
power jets. Different line groups have distinct constant $v_i$;
the parabola and the other isolated group have distinct constant
$u$ values. High-order Hermite cutoffs in these coordinates
separate the groups. Choose their orders to exceed the largest pole order of the
finitely many basis representatives, measured in the common cover.
Then all cross evaluations have positive order. In a smooth
isolated group the same assertion is ordinary finite jet
interpolation. At the parabola chart, the leading denominator
$-2c$ is a unit; multiplying by a high enough power of
$u-c$ separates the other line without changing that unit.
These constructions give an invertible block leading matrix
of exactly the regular rank. There is no cyclic-seed assumption
and no appeal to a count of unconstructed escape characters.

After a common cover this regular part is trivial, with all
original-marker values of positive order: $t^{-1}$ in the
line and parabola models, $t^{-5/6}$ in the cusp's last model,
$t^{-3/4}$ in the tangent's last model, and the positive powers
of the isolated characters. Lemma~\ref{lem:regularpositiveseed}
excludes a constant original marked period. Differential
isolation and shifting $k$ exclude every nonzero finite
exponential-polynomial in the nonisotropic normal forms.

\begin{proposition}\label{prop:cusptangentcomplete}
A pure cubic with cuspidal leading curve or tangent line--conic
leading curve has an isotropic singular direction with $F_2=0$.
It satisfies the mixed cutoff $m>3\deg Q$.
\end{proposition}
\begin{proof}
The gate excludes nonzero $F_2$ at the singular direction.
The exhaustive bounded-group construction excludes the
nonisotropic alternative, including every collision and
degree drop. In the remaining isotropic alternative,
Lemma~\ref{lem:frame} supplies the Gaussian affine-linear
coordinate. The already proved affine-linear theorem gives
the asserted cutoff.
\end{proof}
